\section{Multi-Agent Systems: paralelización, especialización y protocolos de comunicación}

Ante tareas largas, un solo agent está limitado por el contexto, el número de herramientas y la acumulación de errores. Los sistemas multi-agent intentan dividir la tarea entre distintos worker y dejar que un manager coordine los resultados. Así se puede paralelizar y especializar, pero también aparecen nuevos problemas de sistemas distribuidos: la división de tareas, la consistencia de los mensajes, la ejecución duplicada, la espera de sincronización, la propagación de permisos y la recuperación ante fallos.

\subsection{Por qué usar varios Agents}

La primera figura representa los sistemas autónomos con varios robots. El build completo enumera parallelization y task specialization. La paralelización solo reduce la latencia cuando los subtareas dependen poco entre sí. Si todos los worker esperan el mismo estado compartido, el costo de comunicación anula la ganancia.

Esta sección pasa de la memory y la personalization de un solo agent a los límites de la colaboración. Antes de dividir una tarea conviene trazar su grafo de dependencias. Consultar varias fuentes independientes puede hacerse en paralelo, pero redactar el informe final tiene que esperar a que se reúna la evidencia. Cada worker puede usar sus propias herramientas especializadas, pero la identidad del usuario y los permisos de pago, que son compartidos, requieren control central. Al leer la figura no hay que tomar el número de robots como medida de capacidad. Hay que preguntar si cada agent tiene una responsabilidad independiente, si sus salidas pueden combinarse, quién resuelve los conflictos y si el manager aísla a un worker que excede el tiempo límite o devuelve contenido malicioso.

\lecturefigure{slide-52-multi-agent-autonomous-systems.jpg}{Forma básica de un sistema autónomo multi-Agent}{00:39:49}

\lecturefigure{slide-53-why-multiagent-systems.jpg}{Las dos motivaciones principales de los sistemas multi-Agent: paralelización y especialización}{00:40:48}

Si la tarea se divide en $K$ subtareas, la latencia en serie es $\sum_k T_k$. La latencia en paralelo ideal es aproximadamente
\[
T_{\text{parallel}}\approx \max_k T_k+T_{\text{coord}}+T_{\text{merge}}.
\]
$T_{\text{coord}}$ es el tiempo de planificación y comunicación, y $T_{\text{merge}}$ es el de combinación y verificación de resultados. Si el costo de coordinación es demasiado alto o si la subtarea más larga domina, un sistema multi-agent no es más rápido por sí solo.

\begin{warningbox}{«Multi-Agent» no significa copiar el mismo Prompt varias veces}
Un sistema multi-agent sin límites de rol, sin protocolo de estado compartido y sin resolución de conflictos por lo general solo aumenta el costo en tokens. Cada worker debe tener una entrada definida, un schema de salida, permisos, un plazo límite y un manejo de fallos.
\end{warningbox}

\subsection{Hierarchy y syncing primitives}

El diagrama jerárquico entrega la solicitud del usuario a un manager, que luego la reparte entre varios worker. La responsabilidad del manager no se limita a reenviar mensajes: también planifica las dependencias, controla el presupuesto, decide cuando los resultados entran en conflicto y reasigna el trabajo si un worker falla. Las llamadas \term{syncing primitives} son los mecanismos que permiten que varios agent lleguen a un acuerdo operativo sobre el estado compartido, por ejemplo locks, números de versión, barrier, confirmaciones de mensajes e ID de tarea idempotentes.

\lecturefigure{slide-54-agent-hierarchy.jpg}{Jerarquía manager--worker y flujo de información entre Agents}{00:42:21}

\begin{lstlisting}[caption={Contrato de tarea manager--worker}]
task = {
  "task_id": stable_id,
  "goal": scoped_goal,
  "inputs": immutable_inputs,
  "permissions": allowed_tools,
  "depends_on": prerequisite_task_ids,
  "deadline": deadline,
  "output_schema": schema,
  "retry_policy": retry_policy
}
\end{lstlisting}

\subsection{MCP y A2A: un estándar de conexión no es una garantía de seguridad}

La slide de protocolos asigna MCP (Model Context Protocol) a la conexión entre modelos, herramientas y datos, y A2A (Agent-to-Agent) a la comunicación entre agents. El valor de MCP está en unificar la interfaz de descubrimiento e invocación, de modo que cada aplicación no tenga que escribir glue code a la medida para cada herramienta. A2A, en cambio, necesita definir la identidad, los mensajes, el estado de las tareas y la semántica de la colaboración. Ambos resuelven la interoperabilidad, pero no resuelven por sí solos la autorización ni la corrección.

\lecturefigure{slide-55-mcp-a2a-protocols.jpg}{MCP conecta modelos y herramientas; A2A permite la colaboración entre Agents}{00:43:44}

\begin{knowledgebox}{Primer uso: Protocol, Schema y Authorization}
\textbf{Protocol} define cómo se comunican las partes; \textbf{schema} define los campos de los mensajes y sus tipos; \textbf{authorization} determina quién puede ejecutar qué acción. Que un MCP server exponga una herramienta no significa que cualquier client deba tener permiso para invocarla. Los permisos tienen que ser mínimos, contar con capacidad de auditoría y poder revocarse.
\end{knowledgebox}

\subsection{Por qué se compara MCP con USB-C}

La primera figura de MCP destaca la conexión unificada entre modelos, datos y herramientas. El build final enumera además resilience, dynamic tool discovery e intent-based execution. Al leer la figura, «absorber los cambios de la API» debe entenderse como aislamiento mediante una capa de abstracción, no como una interfaz que nunca falla: los cambios en el schema del server, en la semántica y en los permisos siguen requiriendo gestión de versiones y pruebas de compatibilidad.

\lecturefigure{slide-56-mcp-usbc.jpg}{La analogía de MCP con USB-C: conexión unificada entre modelos y herramientas}{00:45:11}

\lecturefigure{slide-57-mcp-api-benefits.jpg}{Resiliencia y descubrimiento dinámico de MCP frente a la integración punto a punto de API}{00:45:46}

La estandarización del protocolo también cambia la localización de fallos. Cuando falla una integración punto a punto, las responsabilidades de la aplicación, del modelo y de la API suelen mezclarse. Con una capa de protocolo se puede revisar por separado si el server es descubrible, si el schema coincide, si los permisos bastan, si la invocación excedió el tiempo límite y si la respuesta pasó la validación. El costo es que la propia capa de protocolo se vuelve una dependencia crítica: hay que versionarla, monitorearla e impedir que un server malicioso inyecte descripciones de herramientas engañosas. La llamada resilience proviene de una separación clara en capas y de las pruebas, no de un nombre unificado.

\teachervoice{00:38:00--00:46:09. El ponente resume el valor de los sistemas multi-agent en parallelization y specialization, y enfatiza hierarchy, synchronization, MCP, A2A, resilient integrations y dynamic discovery. Nota de clase: el lenguaje natural es ambiguo por naturaleza, así que la comunicación necesita un protocolo estructurado y un estado.}

\begin{importantbox}{Lista mínima de seguridad para la capa de protocolo}
El descubrimiento de herramientas debe separarse de la concesión de permisos. Las herramientas sensibles requieren un scope explícito. Cada invocación incluye la identidad de quien llama y el ID de la tarea. Las salidas requieren validación de schema. Las actualizaciones del protocolo deben pasar pruebas de compatibilidad. Las acciones de alto riesgo deben requerir la confirmación del usuario o un control doble.
\end{importantbox}

\subsection{Resumen de la sección}

Los sistemas multi-agent convierten un problema de un solo modelo en un problema de sistemas distribuidos. La paralelización y la especialización solo aportan beneficios cuando la tarea puede dividirse y la comunicación está bajo control. El contrato manager--worker, las primitives de sincronización, los schema de MCP/A2A, la identidad y los permisos determinan en conjunto si el sistema es confiable.
